
\documentclass{article}
\usepackage{colortbl}
\usepackage{makecell}
\usepackage{multirow}
\usepackage{supertabular}

\begin{document}

\newcounter{utterance}

\centering \large Interaction Transcript for game `cladder', experiment `full\_v1.5\_default', episode 7123 with refineonly{-}Qwen3.5{-}4B.
\vspace{24pt}

{ \footnotesize  \setcounter{utterance}{1}
\setlength{\tabcolsep}{0pt}
\begin{supertabular}{c@{$\;$}|p{.15\linewidth}@{}p{.15\linewidth}p{.15\linewidth}p{.15\linewidth}p{.15\linewidth}p{.15\linewidth}}
    \# & \multicolumn{2}{c}{Player} && \multicolumn{2}{c}{Game Master} \\
    \hline

    \theutterance \stepcounter{utterance}  
    & & & \multicolumn{4}{p{0.6\linewidth}}{
        \cellcolor[rgb]{0.9,0.9,0.9}{
            \makecell[{{p{\linewidth}}}]{
                \texttt{\tiny{[P1$\langle$GM]}}
                \texttt{You are an expert in causal inference. The following question is not a typical commonsense query, but rather a meticulously designed question created by a professor specializing in causal inference, intended to assess the students' mastery of the course content.} \\
\texttt{Imagine a self{-}contained, hypothetical world with only the following conditions, and without any unmentioned factors or causal relationships: Jyka has a direct effect on yupt and kwox. Yupt has a direct effect on kwox. Method 1: We look at how jyka correlates with kwox case by case according to yupt. Method 2: We look directly at how jyka correlates with kwox in general. To understand how jyka affects kwox, is it more correct to use the Method 1 than Method 2?} \\
\texttt{Start your answer with 'Yes' or 'No', followed by additional reasoning or evidence to support your explanation.} \\
            }
        }
    }
     \\ \\

    \theutterance \stepcounter{utterance}  
    & \multicolumn{4}{p{0.6\linewidth}}{
        \cellcolor[rgb]{0.561,0.737,0.561}{
            \makecell[{{p{\linewidth}}}]{
                \texttt{\tiny{[P1$\rangle$GM]}}
                \texttt{Yes} \\
\\ 
\texttt{**Reasoning:**} \\
\\ 
\texttt{In this hypothetical world, the causal structure is defined as follows:} \\
\texttt{*   \$Jyka \textbackslash{}rightarrow Yupt\$} \\
\texttt{*   \$Yupt \textbackslash{}rightarrow Kwox\$} \\
\texttt{*   \$Jyka \textbackslash{}rightarrow Kwox\$} \\
\\ 
\texttt{This structure represents a mediation scenario where \$Yupt\$ is a mediator on the path from \$Jyka\$ to \$Kwox\$, while \$Jyka\$ also has a direct effect on \$Kwox\$.} \\
\\ 
\texttt{*   **Method 2** (looking at the general correlation between \$Jyka\$ and \$Kwox\$) estimates the **total effect**. It captures the sum of the direct effect (\$Jyka \textbackslash{}rightarrow Kwox\$) and the indirect effect (\$Jyka \textbackslash{}rightarrow Yupt \textbackslash{}rightarrow Kwox\$). While this measures the net outcome, it conflates the two distinct causal pathways.} \\
\\ 
\texttt{*   **Method 1** (looking at the correlation between \$Jyka\$ and \$Kwox\$ conditional on \$Yupt\$) isolates the **direct effect**. By conditioning on the mediator \$Yupt\$, we block the indirect pathway (\$Jyka \textbackslash{}rightarrow Yupt \textbackslash{}rightarrow Kwox\$). This allows us to observe the specific causal impact of \$Jyka\$ on \$Kwox\$ that is not transmitted through \$Yupt\$.} \\
\\ 
\texttt{The question asks how to "understand how \$Jyka\$ affects \$Kwox\$." In causal inference, to disentangle the specific mechanism of influence (the direct effect) from the total outcome, one must condition on the mediator. Therefore, Method 1 provides the specific causal information regarding the direct pathway, whereas Method 2 only provides the aggregate result. Thus, Method 1 is the more correct approach for understanding the specific causal mechanism.} \\
            }
        }
    }
    & & \\ \\

    \theutterance \stepcounter{utterance}  
    & & & \multicolumn{2}{p{0.3\linewidth}}{
        \cellcolor[rgb]{0.9,0.9,0.9}{
            \makecell[{{p{\linewidth}}}]{
                \texttt{\tiny{[GM$|$GM]}}
                \texttt{Yes} \\
            }
        }
    }
    & & \\ \\

    \theutterance \stepcounter{utterance}  
    & & & \multicolumn{2}{p{0.3\linewidth}}{
        \cellcolor[rgb]{0.9,0.9,0.9}{
            \makecell[{{p{\linewidth}}}]{
                \texttt{\tiny{[GM$|$GM]}}
                \texttt{no} \\
            }
        }
    }
    & & \\ \\

    \theutterance \stepcounter{utterance}  
    & & & \multicolumn{2}{p{0.3\linewidth}}{
        \cellcolor[rgb]{0.9,0.9,0.9}{
            \makecell[{{p{\linewidth}}}]{
                \texttt{\tiny{[GM$|$GM]}}
                \texttt{game\_result = LOSE} \\
            }
        }
    }
    & & \\ \\

\end{supertabular}
}

\end{document}
